\chapter{Outcomes}
\label{chap:bilan}
% Corresponds to phase 04 of the roadmap.

\section{Security review}
A comprehensive, formal security audit (dependency and image scanning, secret rotation policy) is planned for the final phase of the roadmap and had not yet been carried out at the time of writing. What did happen along the way, however, is a series of targeted reviews that already behaved like a lightweight audit and caught several real gaps: an inspection of the live firewall rules revealed that the LAN-only restriction planned for the dashboard and administrative ports had, in fact, never been applied, and that same inspection surfaced an IPv6 exposure that the original CGNAT-based security reasoning had not accounted for; the publication of the dashboard then showed that Docker's own packet-filtering rules bypassed the firewall for the published port; and a check of the effective SSH configuration of the second virtual machine showed it still accepted passwords (see chapter~\ref{chap:securisation}). All were corrected once identified. This is itself a useful finding for the final audit: what a firewall or a configuration is \emph{supposed} to enforce should always be checked against what it \emph{actually} enforces, on every machine and after every change of exposure, rather than assumed from the deployment steps followed.

\section{Results achieved}
By the end of the internship's first roadmap phase, the server-hardening objectives were met in full: application-level authentication with anti-brute-force protection, HTTPS (adapted to the network's CGNAT constraint via Cloudflare Tunnel rather than the originally planned direct Traefik/Let's Encrypt path), a firewall correctly scoped to the local network, and SSH access restricted to key-based authentication with root login disabled. The second phase was also completed: a Jenkins pipeline of four stages, triggered by a webhook, blocks deployment on test failure; backups are checked every night against an actual restore test rather than assumed to be valid; and GitLab runs as a second application on its own virtual machine for the other interns.

The third phase is largely done. Supervision (SSH activity of the virtual machine and of the physical server, container state, Proxmox hosts and disks, disk usage, last Jenkins build, proactive Discord alerts) is in production and was the tool used to validate most of the above in practice. The domain \texttt{ucac-icam.info} was registered and delegated to Cloudflare, and two named tunnels, one per virtual machine, publish the dashboard and GitLab at stable addresses from outside the school network. Ports that Docker had published too widely were closed, and key-only SSH now applies to both virtual machines.

\section{Difficulties and skills gained}
Most of the difficulties encountered during the internship were not solved by knowing the right command in advance, but by a disciplined process of isolating a failure to its actual cause rather than the most visible symptom. A representative example: after fixing the dashboard's login credentials, authentication kept failing even with the correct password. Comparing a request made directly against the VM (bypassing the reverse proxy) with one made through the public URL isolated the problem to the application itself rather than the network path; inspecting the environment variables actually loaded inside the running container (rather than the \texttt{.env} file on disk) revealed the real cause: \texttt{docker compose restart} restarts a container without rereading its environment file, so a credential change only takes effect after \texttt{docker compose up -d --force-recreate}. A second example came with GitLab: after moving it behind the tunnel, login attempts were rejected. The web server's access log showed two different outcomes for the same credentials: a request answered with 200 through the domain, meaning the password itself was refused, and a request answered with 422 on the local address, meaning the session cookie was no longer accepted over plain HTTP. Reading the status codes of the \emph{real} requests, instead of guessing from the error page, separated what looked like a single problem into two distinct ones.

This same method -- narrow the failure down with a controlled comparison, verify the \emph{running} state rather than the \emph{configured} state, read logs before guessing -- reappears throughout chapters~\ref{chap:securisation} and~\ref{chap:cicd}, and is arguably the most transferable skill the internship developed, ahead of any specific tool. Writing each procedure down as soon as a problem was solved had a measurable effect too: the same trap (Docker creating an empty folder in place of a missing file) was documented after its first occurrence and never met a third time.

On the tool side, the internship provided hands-on, production-adjacent experience with virtualisation (Proxmox), containerisation (Docker Compose), reverse proxying and tunnelling, DNS and domain management, CI/CD (Jenkins), the hosting and maintenance of a complex application (GitLab), and Linux system hardening (UFW, SSH, Fail2ban) -- most of it learned by first getting something to fail in a realistic way and then fixing it, rather than by following a guide written for an idealised environment.

\section{Limitations and outlook}
Two open points remain at the time of writing. The first is reliability after an interruption: when the server is switched off and on, the dashboard container does not always come back by itself. The restart policy of the container is correct; the likely cause is the start-up order, because the port is published on a local network address that Docker must find on the network card at the moment it restarts the container, and it can start before the network is configured. The planned fix makes Docker wait for the network to be online and adds a start-up task as a safety net, and it will only be considered done once a real reboot has confirmed it.

The second is protection of what is now public. The two published services still rely on their own authentication only. The planned measures, in order of priority, are: disabling public sign-up on GitLab and enabling two-factor authentication for its administrator; placing an access-control layer (Cloudflare Access) in front of the dashboard, with HTTPS redirection, HSTS and bot protection; enabling DNSSEC for the domain; and replacing the dashboard's direct access to the Docker socket by a filtering proxy, since a read-only mount still allows write calls to the Docker API. Beyond that, hosting a third application would exercise the multi-application mechanism once more, and the formal security audit and the handover documentation close the internship.

\clearpage
